\documentclass[border=4pt]{standalone}
\usepackage{amsmath,amssymb,mathtools,xcolor,varwidth}
\definecolor{radialred}{HTML}{D81B60}
\begin{document}
\begin{varwidth}{28em}
\large\bfseries
The identical reasoning holds at $C$ ($dD \parallel CS$), at $D$, and at every vertex: each deflection lies parallel to the radius, so the impulse always aims at $S$. Taking the polygon to its limit, the \textcolor{radialred}{radial force directions BS, CS, DS} converge continuously on the centre — the force is centripetal. Thus the equal-area observation ALONE locks the direction toward $S$, however the magnitude may vary with distance; combined with Proposition I this makes swept area the definitive test of a central force. \textit{Q.E.D.}
\end{varwidth}
\end{document}
